\chapter{Abenteuer erleben}

Neben Kämpfen erleben Helden viele weitere Situationen. Sie schleichen an Wachen vorbei, erklimmen steile Felswände, überwinden Hindernisse oder unterstützen sich gegenseitig bei schwierigen Aufgaben.

Dieses Kapitel beschreibt die Regeln für Zusammenarbeit und Rast.

\section{Zusammenarbeit}

Helden können Herausforderungen entweder gemeinsam bewältigen oder sich gegenseitig unterstützen.

\subsection{Gruppenproben}

Müssen mehrere Helden gemeinsam dieselbe Herausforderung meistern, handelt es sich um eine Gruppenprobe.

Die Schwierigkeit der Probe entspricht der regulären Schwierigkeit multipliziert mit der Anzahl der beteiligten Helden. Alle beteiligten Helden führen die passende Probe durch und addieren ihre erzielten Sterne. Erreichen oder überschreiten sie gemeinsam die Gesamtschwierigkeit, ist die Gruppenprobe erfolgreich.

\textbf{Beispiel:} Vier Helden möchten sich gemeinsam an einer Wache vorbeischleichen. Die Schwierigkeit beträgt 4.

Da vier Helden beteiligt sind, beträgt die Gesamtschwierigkeit 16 Sterne. Alle vier Helden führen eine Schleichen-Probe durch. Erzielen sie gemeinsam mindestens 16 Sterne, bleibt die Gruppe unentdeckt.

\subsection{Helfen}

Ein Held kann einen anderen bei einer Probe unterstützen.

Dazu beschreibt der Spieler, wie sein Held helfen möchte. Ist die Unterstützung sinnvoll, führt der unterstützende Held zuerst eine passende Probe durch. Welche Fertigkeit dafür verwendet wird, entscheiden Spieler und Spielleitung gemeinsam.

Für jeden erzielten Stern erhält der aktive Held einen zusätzlichen blauen Würfel für seine Probe.

\section{Rast}

Helden können rasten, wenn sie ausreichend Zeit und Sicherheit dafür haben. Ob und wann eine Rast möglich ist, entscheidet die Spielleitung.

Jeder Held rastet unabhängig von den anderen. Während einige Helden rasten, können andere beispielsweise Wache halten oder anderen Tätigkeiten nachgehen.

Es zählen nur vollständig abgeschlossene Stunden einer Rast.

Für jede vollständige Stunde Rast erhält ein Held:

\begin{itemize}
  \item Lebenspunkte in Höhe seines KON-Wertes, höchstens bis zu seinem Maximum.
  \item Mana in Höhe seines CHA-Wertes, höchstens bis zu seinem Maximum.
\end{itemize}

Nach einer vollständig abgeschlossenen Stunde Rast werden außerdem alle umgedrehten Karten -- beispielsweise Talente, Zauber und Ausrüstung -- wieder aufgedeckt.